\clearpage
\section{Basketball}
\subsection{random variables}
\emph{Who's Shimshon? I'm assuming this is supposed to be Gal :)}

We can define \(S_1,S_2,S_3\) to be the score achieved through penalties, inside, and outside throws. This means the total score he earned is
\begin{equation}
    S = S_1+S_2+S_3
\end{equation}
but the number of successful throws is distributed binomial so 
\begin{equation}
    S_1\sim Bin\ins{N_1,q_1},\quad S_2\sim Bin\ins{N_2,q_2},\quad S_3\sim Bin\ins{N_3,q_3}
\end{equation}
noting the large Ns are different from small ns. N represents an unknown number of attempts for the distribution, n is the number actually attempted.
\subsection{Average score}
\emph{Who's Shimshon? I'm assuming this is supposed to be Gal :)}

\begin{equation}
    E[S]=E[S_1+2S_2+3S_3]=E[S_1]+2E[S_2]+3E[S_3]
\end{equation}
but 
\begin{align}
    E[S_1]=n_1q_1=6\cdot0.75=4.5,\qquad E[S_2]=6,\qquad E[S_3]=4
\end{align}
so 
\begin{equation}
    E[S]=28.5
\end{equation}
and the variance is, since each set is independent of the others
\begin{align}
    V[S]&=V[S_1+2S_2+3S_3]=V[S_1]+4V[S_2]+9V[S_3]=n_1q_1\ins{1-q_1}+4n_2q_2\ins{1-q_2}+9n_3q_3\ins{1-q_3}\\
    &= 19.08 \implies \sigma= 4.37
\end{align}
\subsection{Most probable throws}
The assumption \(H_0\) implies 
\begin{equation}
    0.75S_1+2\cdot0.75S_2+3\cdot0.66S_3=24
\end{equation}
together with 
\begin{equation}
    0.25S_1+0.25S_2+0.34S_3=5
\end{equation}
and 
\begin{equation}
    S_1+S_2+S_3=19
\end{equation}
which is three equations for three variables so 
\begin{equation}
    \mqty(1/2&3/2&2\\1/4&1/4&1/3\\1&1&1)\mqty(S_1\\S_2\\S_3)=\mqty(24\\5\\19)
\end{equation}
throwing into wolfram alpha we get 
\begin{equation}
    S_1=8,\qquad S_2=8,\qquad S_3=3
\end{equation}
\subsection{Extension}
\emph{Who's Shimshon? I'm assuming this is supposed to be Gal :)}

after the extension the numbers are 
\begin{equation}
    S_1=10,\quad S_2=8,\quad S_3=6
\end{equation}
so the p-value is 
\begin{align}
    &P(S_1<8\cdot0.75,S_2<8\cdot0.75,S_3<3\cdot0.66|H_0)\\
    &=P(S_1<6,S_2<6,S_3<2|H_0)\\
    &= P(S_1<6)P(S_2<6)P(S_3<2)
\end{align}
which can be directly computed using binomial distribution probability distributions.